\section{Introduction}
\label{sec:introduction}

Many real-world engineering and scientific problems require the optimisation of
complex, high-dimensional objective functions for which gradient information is
unavailable, unreliable, or prohibitively expensive to compute.  Classical
mathematical programming methods---such as Newton's method, conjugate gradient
descent, or interior-point algorithms---rely on smoothness, convexity, or at
least local differentiability assumptions that are routinely violated in
practice.  Combinatorial landscapes, black-box simulations, and functions
riddled with local optima demand a fundamentally different class of solvers.
Metaheuristic algorithms address this need by conducting stochastic, often
population-based searches that sacrifice guarantees of global convergence in
favour of robust, general-purpose exploration of the decision space.  Over the
past three decades, the field has produced a rich taxonomy of such methods,
including genetic algorithms, differential evolution, ant colony optimisation,
and, most pertinently to the present work, particle swarm optimisation (PSO).

PSO was introduced by Kennedy and Eberhart in 1995 as a socially inspired
optimisation paradigm in which a population of candidate solutions---called
\emph{particles}---traverse the search space under the influence of their own
historical best positions and the best position discovered by the swarm as a
whole~\cite{Kennedy1995}.  Each particle maintains a velocity vector that is
updated at every iteration according to cognitive and social acceleration
coefficients, producing a balance between individual memory and collective
intelligence.  The elegance and simplicity of the original formulation
attracted wide attention, and subsequent refinements strengthened its
theoretical foundations and practical performance.  Shi and Eberhart introduced
the \emph{inertia weight}, a damping factor applied to the previous velocity
that enables a smooth transition from exploration to exploitation as the
optimisation progresses~\cite{Shi1998}.  Clerc and Kennedy proposed the
\emph{constriction coefficient} framework, derived from a stability analysis of
the swarm dynamics, which automatically constrains the velocity magnitudes and
guarantees convergence to a fixed point under idealised
conditions~\cite{Clerc2002}.  Together, these contributions established PSO as
a mature and competitive optimiser across a broad range of continuous
optimisation benchmarks.

Despite its success, classical PSO exhibits well-documented limitations.  The
velocity-based update mechanism can lead to premature convergence when the
inertia weight decays too aggressively, or to oscillatory, non-convergent
behaviour when the acceleration coefficients are poorly tuned.  On highly
multimodal landscapes such as the Rastrigin or Schwefel functions, the swarm
frequently collapses into a single basin of attraction, foregoing large regions
of the search space that may contain the global optimum.  These difficulties
have motivated the exploration of alternative update mechanisms that draw on
principles from quantum mechanics.

Sun, Feng, and Xu proposed Quantum-Behaved Particle Swarm Optimisation (QPSO)
in 2004 as a response to these shortcomings~\cite{Sun2004}.  In QPSO,
particles no longer carry velocity vectors.  Instead, their positions are
sampled from a probability distribution derived from a quantum-mechanical
potential well centred on an \emph{attractor point}---a stochastic
interpolation between each particle's personal best and the swarm's global
best.  The width of this distribution is governed by a contraction--expansion
coefficient~$\beta$ and by the distance between the particle's current
position and the \emph{mean best position} (mbest), defined as the centroid of
all personal bests.  This formulation endows the swarm with a quantum-inspired
tunnelling capability: even when a particle is close to a local minimum, there
exists a non-zero probability that it will appear on the far side of a
potential barrier in a single iteration, a phenomenon with no analogue in the
deterministic velocity updates of classical PSO.  Sun et al.\ later provided a
rigorous analysis of individual particle behaviour in QPSO, establishing
convergence conditions and parameter selection guidelines~\cite{Sun2012}.

The randomness that drives both PSO and QPSO---the vectors $r_1$ and $r_2$ in
classical PSO, and the attractor interpolation, displacement magnitude, and
tunnelling sign in QPSO---is conventionally generated by pseudo-random number
generators (PRNGs).  Although modern PRNGs produce sequences that pass
stringent statistical tests for uniformity and independence, they are
inherently deterministic: given the same seed, the same trajectory is
reproduced exactly.  This raises a natural question: \emph{does replacing
pseudo-random numbers with genuinely quantum-random numbers alter the search
dynamics or solution quality of QPSO?}  The question is not merely academic.
Quantum random number generators (QRNGs), now accessible through cloud-based
quantum computing platforms such as IBM Qiskit, produce bit strings via
physical processes---superposition and measurement collapse---that are
fundamentally unpredictable according to the axioms of quantum mechanics.  If
the statistical structure of the random inputs matters to the optimiser's
trajectory, then quantum randomness could, in principle, confer an advantage.

Surprisingly, despite the growing accessibility of quantum hardware simulators
and the intense interest in quantum-enhanced algorithms, no systematic
empirical study has compared QPSO under different randomness
sources---classical PRNG, fully quantum, and hybrid---using a unified
implementation and a rigorous statistical framework.  Prior work has either
focused on the mathematical formulation of QPSO without varying the randomness
source, or has employed quantum circuits for specific sub-tasks without
controlling for the remaining sources of stochasticity.  This gap leaves
practitioners without evidence-based guidance on whether the computational
overhead of quantum circuit execution is justified by measurable improvements
in optimisation outcomes.

The present work addresses this gap through three principal contributions.
First, we develop a \emph{unified software framework} that implements both
classical PSO and QPSO within a single, modular codebase and exposes three
randomness modes for QPSO: (i)~\texttt{math} mode, in which all random values
are drawn from NumPy's Mersenne Twister PRNG; (ii)~\texttt{full} mode, in
which every random quantity---cognitive and social weights, displacement
magnitudes, and tunnelling signs---is generated by Qiskit quantum circuits
comprising Hadamard gates and parameterised $R_Y$ rotations; and
(iii)~\texttt{hybrid} mode, in which only the binary tunnelling sign is drawn
from a quantum circuit while all continuous random variables remain classical.
The framework incorporates established enhancements from the literature,
including linear inertia decay and velocity clamping for
PSO~\cite{Shi1998,Clerc2002}, and the weighted mean best position of Xi, Sun,
and Xu for QPSO~\cite{Xi2008}.

Second, we conduct a \emph{comprehensive benchmark study} across six standard
test functions---Sphere, Rastrigin, Rosenbrock, Ackley, Griewank, and
Schwefel---spanning unimodal, multimodal, valley-shaped, and deceptive
landscape topologies, at dimensionalities $n \in \{2, 5, 10, 20\}$.  Classical
PSO and QPSO (in the mathematical randomness mode) are each evaluated over 50
independent runs with controlled, paired random seeds, and we record
convergence histories, iteration-to-success counts, and final objective
values.

Third, we perform a \emph{rigorous statistical comparison} of the results,
employing the non-parametric Wilcoxon signed-rank test---complemented, for the
multimodal Rastrigin function, by a binary local-minimum escape-rate analysis
under Fisher's exact test---to determine whether observed differences between
the two algorithms are statistically significant and practically meaningful,
rather than artefacts of sampling variability.

The remainder of this paper is organised as follows.
Section~\ref{sec:background} reviews the mathematical foundations of classical
PSO and QPSO, the weighted mean best extension, and the benchmark functions
used in our experiments.  Section~\ref{sec:methodology} describes the software
architecture, the three quantum randomness modes, and the convergence
mechanisms employed.  Section~\ref{sec:results} presents the experimental
results, statistical analysis, and discusses their implications.  Finally,
Section~\ref{sec:conclusion} summarises the contributions and suggests
directions for future research.
